% FastaGuard — Bioinformatics (OUP) Application Note draft
% Target: Application Note (~4 pages / ~2,000–2,600 words)
% Initial submissions may be format-free; use OUP General Template for revision.
% See README.md for build steps and venue notes.
%
% Compile: make   (or: latexmk -pdf main.tex)
\pdfminorversion=5
\documentclass[11pt,a4paper]{article}

% Draft layout approximating journal constraints without requiring oup-authoring-template.
% Switch to OUP General Template (\documentclass[numsec,webpdf,modern,large]{...})
% before camera-ready / revision formatting.
\usepackage[margin=2.5cm]{geometry}
\usepackage{times}
\usepackage{graphicx}
\usepackage{booktabs}
\usepackage{hyperref}
\usepackage{xcolor}
\usepackage{natbib}
\usepackage{url}
\usepackage{setspace}
\usepackage{lineno}

\hypersetup{
  colorlinks=true,
  linkcolor=black,
  citecolor=blue!50!black,
  urlcolor=blue!50!black
}

% Double-spacing + line numbers help page-budget checks for first submission.
\doublespacing
\linenumbers

\title{FastaGuard: FASTA preflight QC for modern bioinformatics pipelines}

\author{
  Ehsan Estaji%
  \thanks{Corresponding author. Email: \texttt{[institutional-email@example.org]}}%
  \\
  \textit{[Affiliation to be completed]}
}

\date{}

\begin{document}
\maketitle

\begin{abstract}
\noindent\textbf{Summary:}
FastaGuard is a fast, explainable FASTA preflight QC tool that validates assembly
FASTA files, detects structural and composition red flags, and emits a stable,
pipeline-ready report contract (HTML, JSON, TSV, MultiQC-compatible JSON) before
expensive downstream steps such as QUAST, BUSCO, BlobToolKit, CheckM, or annotation.

\noindent\textbf{Availability and Implementation:}
Source, binaries, and documentation are freely available under the MIT license at
\url{https://github.com/ehsanestaji/FastaGuard}.
Install via Bioconda (\texttt{fastaguard}), BioContainers, or GitHub release binaries.
Implemented in Rust; supported on Linux and macOS.

\noindent\textbf{Contact:}
\texttt{[institutional-email@example.org]}

\noindent\textbf{Supplementary information:}
Supplementary figures, finding catalog excerpts, and reproducibility notes will
accompany the submitted manuscript (see \texttt{paper/} in the repository).
\end{abstract}

\section*{Motivation}
FASTA files underpin genome assembly workflows, yet FASTA-level quality control
remains fragmented across ad hoc scripts, general sequence utilities, and heavy
downstream evaluators~\citep{gurevich2013quast,simao2015busco,challis2020blobtoolkit}.
Malformed structure, duplicate identifiers, invalid alphabet characters, high-N
content, and submission hazards often surface only after costly QC or official
repository validation has already started.

FastaGuard fills the missing \emph{preflight} layer: a single, opinionated first
command that answers whether a FASTA is valid, sane, interpretable, and safe to
pass downstream, and that returns a deterministic PASS/WARN/FAIL contract for
humans, workflow engines, and tool agents~\citep{ewels2016multiqc}.
It is not a replacement for interpretive assembly QC; it is the earlier gate
that makes those tools safer and cheaper to run.

\section*{Implementation and features}
FastaGuard is a streaming-first Rust CLI focused on assembly FASTA preflight.
Core capabilities include:

\begin{itemize}
  \item Structural validity checks (headers, empty records, parse failures) with
        explainable FAIL verdicts.
  \item Identifier and alphabet hygiene (duplicate IDs/sequences, invalid IUPAC).
  \item Assembly stats and composition signals (N50/L50 family metrics; GC/AT/N
        and ambiguity rates; gap runs; tiny-contig advisories; outlier hints).
  \item Pipeline gates (\texttt{--gate pipeline}) exposing
        \texttt{gate.can\_continue}, blocking vs advisory findings, and input
        checksum provenance for stop/go routing in Nextflow, Snakemake, nf-core,
        and related engines.
  \item Submission-readiness preflight (\texttt{--gate submission}) for FASTA-level
        NCBI genome policy signals without replacing official validators.
  \item Cohort triage via \texttt{fastaguard compare} across many assemblies.
  \item Reference compatibility mode for FASTA vs index/dictionary/header
        declaration agreement (v1 reference contract).
  \item Stable multi-format outputs: HTML report, versioned JSON schema, TSV,
        and MultiQC-compatible JSON.
\end{itemize}

\noindent
Install and smoke-test paths are intentionally low-friction (Bioconda,
BioContainers, and release binaries), so Application Note readers can reproduce
basic runs without heavy local builds.

% Placeholder figure — replace with real pipeline schematic or benchmark plot.
\begin{figure}[t]
  \centering
  \fbox{\parbox{0.92\linewidth}{\centering\vspace{2.2cm}
  \textbf{[Figure~1 placeholder]}\\[0.4em]
  FastaGuard as the preflight layer before QUAST / BUSCO / BlobToolKit / MultiQC.\\[0.3em]
  \textit{Replace with \texttt{figures/preflight-overview.pdf} or PNG.}
  \vspace{2.2cm}}}
  \caption{Conceptual placement of FastaGuard relative to common assembly QC tools.
  FastaGuard answers FASTA validity and triage questions; interpretive QC remains
  downstream.}
  \label{fig:overview}
\end{figure}

\section*{Results and use}
On representative public and synthetic assemblies, FastaGuard completes preflight
in seconds-scale wall time with modest memory, returning machine-readable
verdicts suitable for workflow branching
(see Table~\ref{tab:benchmark-placeholder} and repository evidence under
\texttt{docs/evidence/}).
Typical outcomes are: PASS with empty finding sets on clean references;
non-blocking WARN advisories (for example gap-pattern hints) that route follow-up
QC; and FAIL when structure, alphabet, or gate-blocking conditions are met.

% Placeholder table — replace with measured public-evidence numbers before submission.
\begin{table}[t]
  \centering
  \caption{Illustrative preflight outcomes (replace with frozen public-evidence
  numbers before submission).}
  \label{tab:benchmark-placeholder}
  \begin{tabular}{@{}lllrr@{}}
    \toprule
    Input role & Verdict & Gate continue & Time & Peak RSS \\
    \midrule
    Clean public assembly & PASS & yes & TBD s & TBD MB \\
    Public assembly + advisories & WARN & yes & TBD s & TBD MB \\
    Synthetic contract FAIL case & FAIL & no & TBD s & TBD MB \\
    \bottomrule
  \end{tabular}
\end{table}

FastaGuard is already packaged for community workflows (Bioconda, BioContainers,
merged nf-core module and Snakemake wrapper updates). The citable artifact is
the versioned software contract plus this Application Note.

\section*{Discussion}
Positioning FastaGuard as preflight—not as FastQC-for-FASTA or as a QUAST/BUSCO
competitor—keeps claims falsifiable and useful: catch FASTA-level problems early,
emit a stable schema, and let interpretive tools do interpretive work.
Remaining work for the published note includes locking public-evidence tables,
one overview figure, and a short finding-catalog supplement.

\section*{Acknowledgements}
% Add acknowledgements, funding, and conflict disclosures as required by the venue.
The author thanks early adopters and workflow maintainers who exercised packaging
and module integrations. Funding: [to be completed].

\bibliographystyle{plainnat}
\bibliography{references}

\end{document}
